% META: {"id":"1SPE-DERLOCAL-EX-043","chapitre":"1SPE-DERIVATION-LOCAL","type_objet":"exercice","capacites":["1SPE-DERIVATION-LOCAL-C3"],"capacites_codes":["C3"],"methodes":["M1"],"parcours":2,"competences":["calculer","representer"],"duree_min":15,"mode_creation":"ex_nihilo","sources_inspiration":[],"fichier_tex":"chapitres/1SPE-DERIVATION-LOCAL/exercices/1SPE-DERLOCAL-EX-043.tex","corrige_tex":"chapitres/1SPE-DERIVATION-LOCAL/corriges/1SPE-DERLOCAL-CO-043.tex","status":"approved","origin":"generated","status_history":["generated","audited","accepted","approved"],"release_acceptance":"RELEASE_OWNER_BATCH_ACCEPTANCE","acceptance_closure_digest":"sha256:9b3ccf9a81c5520fbb7e03b7b2d3e7bf2057a02834b6d908bf3be5f49f3b553f","acceptance_receipt_digest":"sha256:4fad86f0282e0fa4e0419cca1ad6379ae7af5a280c04a4a90880f90a1c044242"}
% BEGIN-VERIFY
% from sympy import symbols
% x = symbols('x')
% f = -(x - 2)**2 + 5
% fp = f.diff(x)
% assert f.subs(x, 0) == 1
% assert fp.subs(x, 0) == 4
% assert f.subs(x, 2) == 5
% assert fp.subs(x, 2) == 0
% assert f.subs(x, 4) == 1
% assert fp.subs(x, 4) == -4
% END-VERIFY
\begin{exercice}{1SPE-DERLOCAL-EX-043}{2}{15}
On considère la fonction $f$ définie sur $\mathbb{R}$ par $f(x) = -(x-2)^2 + 5$.

On admet que $f$ est dérivable sur $\mathbb{R}$ et que $f'(x) = -2x + 4$.

La courbe $\mathcal{C}_f$ passe par les points $A(0\,;\,1)$, $B(2\,;\,5)$ et $C(4\,;\,1)$.
\begin{enumerate}
\item Calculer $f'(0)$, $f'(2)$ et $f'(4)$.
\item Que représente géométriquement le nombre $f'(2) = 0$ pour la courbe $\mathcal{C}_f$ ?
\item En déduire l'équation de la tangente à $\mathcal{C}_f$ au point $B$.
\item Le point $B$ est-il un maximum local de $f$ ? Justifier à l'aide du signe de $f'$.
\end{enumerate}
\end{exercice}
